% Generated draft: requires mathematical and source review.
\documentclass[11pt]{article}
\usepackage[margin=25mm]{geometry}
\usepackage{fontspec}
\setmainfont{Latin Modern Roman}
\usepackage{amsmath,amssymb,mathtools,unicode-math}
\setmathfont{Latin Modern Math}
\usepackage{enumitem,needspace,xurl,hyperref}
\setlength{\parindent}{0pt}
\setlength{\parskip}{5pt}
\setlength{\emergencystretch}{3em}
\newcommand{\N}{\mathbb{N}}
\newcommand{\Z}{\mathbb{Z}}
\newcommand{\Q}{\mathbb{Q}}
\newcommand{\R}{\mathbb{R}}
\newcommand{\C}{\mathbb{C}}
\newcommand{\F}{\mathbb{F}}
\newcommand{\A}{\mathbb{A}}
\newcommand{\PP}{\mathbb{P}}
\newcommand{\E}{\mathbb{E}}
\newcommand{\eps}{\varepsilon}
\newcommand{\dd}{\,\mathrm{d}}
\begin{document}
\section*{UnsolvedMath ID: 800675}
\textbf{Source problem code:} catalogue entry
\subsection{Definitions and mathematical statement}
Let \(\mathscr F\) be the holomorphic foliation in the supplied source statement, of degree \(d\geq 2\), with non-hyperbolic singularities. Write \(X\) for its ambient complex manifold and \(S=\operatorname{Sing}(\mathscr F)\). The supplied excerpt does not specify \(X\), the definition of degree, or the convention for non-hyperbolic singularities; these are unresolved parts of the hypotheses. In particular, it does not explicitly say whether every singularity must be non-hyperbolic or whether the presence of such singularities suffices.

In the usual setting of a one-dimensional holomorphic foliation, a positive current \(T\) of bidimension \((1,1)\) is directed by \(\mathscr F\) if, on \(X\setminus S\),
\[
T\wedge\alpha=0
\]
for every local holomorphic \(1\)-form \(\alpha\) annihilating the tangent direction of \(\mathscr F\). Positivity means nonnegative pairing with compactly supported smooth strongly positive test forms of complementary bidegree. For definiteness, take \(d^c=\frac{i}{2}(\bar\partial-\partial)\), so that \(dd^c=i\partial\bar\partial\); the vanishing condition is independent of a nonzero normalization factor. A current is \(dd^c\)-closed if \(dd^cT=0\), and is closed if \(dT=0\), in the sense of distributions.

A usual meaning of diffuse for a directed positive current is that its trace measure gives zero mass to each individual regular leaf, equivalently that the corresponding transverse measures have no atoms where such a representation is available. The supplied excerpt does not define diffuse or prescribe behavior on \(S\), so the intended convention requires confirmation.

\textbf{Problem.} Under the source's intended ambient and singularity hypotheses, describe all positive \(dd^c\)-closed currents directed by \(\mathscr F\). Determine under which conditions there exists a nonzero diffuse positive closed current directed by \(\mathscr F\). The ambient setting, the singularity quantifier, and the convention for diffuse currents must be recovered from the source before this becomes a fully specified classification and existence problem.
\subsection{Short English statement}
For a holomorphic foliation of degree \(d\geq 2\) with non-hyperbolic singularities, describe its directed positive \(dd^c\)-closed currents and determine when nonzero diffuse directed positive closed currents exist. The supplied excerpt leaves the ambient space and several conventions unspecified.
\subsection{Sources}
\begin{itemize}
\item Local UnsolvedMath registry: permanent ID \(800675\), exact source problem code \texttt{catalogue entry}, title ``Holomorphic Foliation Theory --- Problem 2,'' in the AMR Open Problem Lists collection. The registry attributes the problem to Dinh and Sibony, with proposed year 2018, and labels it \texttt{open}; this is a supplied status label, not an independently verified literature assessment.
\item Dinh, Sibony, \emph{Some Open Problems on Holomorphic Foliation Theory} (2018), Problem 2. Supplied source URL: \url{https://arxiv.org/abs/1807.05640}. The supplied provenance records access on 2026-10-16 and extraction by \texttt{arxiv\_latex}; these provenance claims have not been independently verified here.
\end{itemize}
\end{document}
